\section{Self-Consistent Fixed Point}

\subsection{Effective bridge capacity}

The previous sections established the base electromagnetic bridge count
\[
B=137
\]
and the net first-order feedback burden
\[
K=\frac{74}{15}.
\]
The inverse compact-fiber electromagnetic bridge coupling is denoted
\[
x=\alpha_F^{-1}.
\]
The quantity \(x\) is not a primitive carrier cardinality. The primitive material carrier has
\[
|F|=128,
\]
and the base bridge count is
\[
B=128+9=137.
\]
The effective bridge capacity \(x\) is obtained only after cover-mediated feedback is admitted and distributed over the same bridge capacity that it modifies.

The central point is that \(K\) is not a new carrier pool and not a direct count to be added to \(B\). It is a feedback burden delivered into the electromagnetic bridge. The appropriate first-order equation is therefore self-consistent:
\[
x=B+\frac{K}{x}.
\]

\subsection{\texorpdfstring{Why the feedback term is \(K/x\)}{Why the feedback term is K/x}}

If the feedback burden \(K\) were an independent carrier pool, the first-order equation would be
\[
x=B+K.
\]
That is not its role. The burden is delivered into the existing bridge and distributed over the effective bridge capacity. A burden distributed over \(N\) effective channels contributes as
\[
\frac{K}{N}.
\]

If the bridge capacity were fixed at its uncorrected base value \(B\), the one-pass perturbative correction would be
\[
x\approx B+\frac{K}{B}.
\]
But the feedback changes the bridge capacity that appears in the denominator. The denominator must therefore be the final effective capacity
\[
x,
\]
not the uncorrected base count
\[
B.
\]
Thus the first-order feedback contribution is
\[
\frac{K}{x}.
\]

\subsection{First-order bridge equation}

The first-order compact-fiber bridge equation is
\[
\boxed{
x=B+\frac{K}{x}.
}
\]
Substituting
\[
B=137,
\qquad
K=\frac{74}{15},
\]
gives
\[
x=137+\frac{74}{15x}.
\]
Multiplying by \(x\),
\[
x^2=137x+\frac{74}{15}.
\]
Therefore
\[
\boxed{
x^2-137x-\frac{74}{15}=0.
}
\]
This is the first-order fixed-point equation for the effective electromagnetic bridge capacity.

\subsection{Positive root}

The positive root is
\[
x_{\rm quad}
=
\frac{137+\sqrt{137^2+4(74/15)}}{2}.
\]
Since
\[
4\left(\frac{74}{15}\right)
=
\frac{296}{15},
\]
this may also be written as
\[
x_{\rm quad}
=
\frac{137+\sqrt{137^2+296/15}}{2}.
\]
Numerically,
\[
\boxed{
x_{\rm quad}=137.03600027236\ldots.
}
\]
The corresponding first-order compact-fiber bridge coupling is
\[
\alpha_{F,\rm quad}
=
\frac{1}{x_{\rm quad}}.
\]

\subsection{First-order accuracy and residual scale}

The first-order result is structurally nontrivial. It fixes the inverse electromagnetic bridge coupling close to the observed low-energy inverse fine-structure constant scale without using the measured value of \(\alpha\) as an input.

It is not, however, the final compact-fiber prediction. At modern precision, the first-order value remains high by approximately
\[
8\text{ ppb}
\]
relative to the quoted low-energy inverse fine-structure constant. The remaining discrepancy motivates the closed-return analysis of the admitted first-order feedback burden.

\subsection{Independence from the closed-return correction}

The first-order equation
\[
x=B+\frac{K}{x}
\]
depends only on the self-consistent transport-sharing rule: a feedback burden delivered into an effective bridge capacity is distributed over the final capacity that it helps determine.

Under that rule, the first-order result follows algebraically. It does not depend on the second-order closed-return correction. If the closed-return term were revised, the first-order compact-fiber bridge lemma would remain intact.